\documentclass[11pt,a4paper]{article}
\usepackage[margin=25mm]{geometry}
\usepackage{fontspec}
\setmainfont{TeX Gyre Pagella}
\usepackage{amsmath,amssymb,amsthm,mathtools}
\usepackage{enumitem,microtype,xcolor}
\usepackage[colorlinks=true,linkcolor=blue!45!black,citecolor=blue!45!black,urlcolor=blue!45!black]{hyperref}
\usepackage{fancyhdr}
\pagestyle{fancy}
\fancyhf{}
\fancyhead[L]{\small Nonboundary localization quotients}
\fancyhead[R]{\small Research note}
\fancyfoot[C]{\thepage}
\setlength{\headheight}{14pt}
\setlength{\emergencystretch}{3em}
\setlength{\parskip}{3pt}
\setlist{itemsep=3pt,topsep=5pt}
\allowdisplaybreaks[2]
\newcommand{\C}{\mathbb C}
\newcommand{\Z}{\mathbb Z}
\newcommand{\g}{\mathfrak g}
\newcommand{\slc}{\mathfrak{sl}_2}
\newcommand{\ad}{\operatorname{ad}}
\newcommand{\Ann}{\operatorname{Ann}}
\newcommand{\Hom}{\operatorname{Hom}}
\newcommand{\Span}{\operatorname{span}}
\newcommand{\Supp}{\operatorname{Supp}}
\newcommand{\pole}{\operatorname{pole}}
\newcommand{\I}{\mathcal I}
\newcommand{\X}{\mathcal X}
\newtheoremstyle{plainupright}{6pt}{5pt}{\normalfont}{}{\bfseries}{.}{.5em}{}
\theoremstyle{plainupright}
\newtheorem{proposition}{Proposition}[section]
\newtheorem{lemma}[proposition]{Lemma}
\newtheorem{corollary}[proposition]{Corollary}
\newtheorem{remark}[proposition]{Remark}
\hypersetup{pdftitle={Nonboundary Localization Quotients of Smooth Affine sl2 Modules},
 pdfsubject={Definitions, proofs, and open problems for quotient localizations}}
\title{\vspace{-14mm}Nonboundary Localization Quotients\\
of Smooth Affine \(\slc\)-Modules\\[3pt]
\large The basic module \(T_{f_0}M_{-1,1}\) and the family \(T_{f_c}M_{a,b}\), \(c<b\)}
\author{}
\date{September 11, 2026}
\begin{document}
\maketitle
\vspace{-9mm}
\begin{abstract}
We collect the currently established structural properties of nonboundary
localization quotients. These include PBW bases, exact root-mode and positive
Cartan profiles, cyclicity, an explicit proper submodule at integral rank-one
parameters in the deep region, annihilator preservation, and compatibility
with further localization and adjoining the degree derivation.
The basic example \(T_{f_0}M_{-1,1}\) is treated explicitly.
Simplicity at nonintegral parameters and a full submodule classification remain
open here. The propositions below have direct proofs; this note is not a claim
of independent peer review or formal verification.
\end{abstract}

\section{Definitions, conventions, and hypotheses}

All vector spaces and tensor products are over \(\C\). We initially work
without the degree derivation:
\[
\g=(\slc\otimes\C[t,t^{-1}])\oplus\C K,\qquad U=U(\g).
\]
Write \(e_r=e\otimes t^r\), \(f_r=f\otimes t^r\), and \(h_r=h\otimes t^r\).
Our normalization is
\begin{align*}
[e_r,f_s]&=h_{r+s}+r\delta_{r+s,0}K,&
[h_r,h_s]&=2r\delta_{r+s,0}K,\\
[h_r,e_s]&=2e_{r+s},&
[h_r,f_s]&=-2f_{r+s}.
\end{align*}
The \(e\)'s commute with one another, the \(f\)'s commute with one another,
and \(K\) is central. Here \(\delta\) denotes the Kronecker delta.

Fix \(a,b\in\Z\) with \(a+b\ge0\), and set \(L=a+b+1\ge1\). Let
\[
S_{a,b}=\Span\{K,\ h_i\ (i\ge0),\ e_i\ (i>a),\ f_j\ (j>b)\}.
\]
Fix the character \(\chi:S_{a,b}\to\C\) specified by
\[
\chi(e_i)=\chi(f_j)=0,\qquad
\chi(K)=\kappa,\qquad \chi(h_i)=\lambda_i,\qquad
\lambda_i=0\quad(i>L).
\]
These conditions make \(\chi\) a Lie algebra character.
Let \(\C_\chi\) be the one-dimensional \(S_{a,b}\)-module defined by \(\chi\), and put
\[
M=M_{a,b}(\chi)=U\otimes_{U(S_{a,b})}\C_\chi,\qquad u=1\otimes1_\chi.
\]
The notation \(\kappa\) distinguishes the central parameter from the
localization index \(c\).

Fix an integer \(c<b\), and write
\[
F=f_c,\qquad d_*=b-c\ge1,\qquad
U_F=U[F^{-1}],\qquad D_FM=U_F\otimes_U M.
\]
The powers of \(F\) form an Ore set, since \(\ad F\) is locally nilpotent on
\(U\). PBW shows that \(F\) is injective on \(M\), so we identify \(M\)
with its image in \(D_FM\) and define
\[
Q_c=T_FM=D_FM/M,\qquad w_m=F^{-m}u+M\quad(m\ge1).
\]
Negative powers are interpreted in \(D_FM\).
\emph{There is no operator \(F^{-1}\) on \(Q_c\).}

A module is \emph{smooth} if each vector is annihilated by
\(\slc\otimes t^N\C[t]\) for some \(N\).
An operator \(x\) is locally nilpotent if each vector is killed by a power of
\(x\); it is locally finite if each vector spans a finite-dimensional
\(\C[x]\)-module. We call \(v\) \(F\)-torsion if \(F^n v=0\) for some \(n\).
Weights below refer to \(\C h_0\oplus\C K\), not to all \(h_i\).

Unless explicitly stated otherwise, no assumption is made on
\(\lambda_L\) or \(\kappa\). For comparison, the original induced module
satisfies \cite[Theorem 1.1]{FGXZ}
\begin{equation}\label{eq:original-simple}
M_{a,b}(\chi)\text{ is simple}\quad\Longleftrightarrow\quad
\lambda_L\ne0\ \text{ and }\ \kappa\ne-2.
\end{equation}
This criterion cannot be transferred directly to \(Q_c\).

\section{PBW structure, root modes, and weights}

\begin{proposition}[PBW model and pole filtration]\label{prop:pbw}
Choose a PBW order with \(F\) first, followed by the families
\[
f_j\ (j\le b,\ j\ne c),\qquad e_i\ (i\le a),\qquad h_r\ (r<0),
\]
and finally a basis of \(S_{a,b}\). Let \(\X\) denote the ordered monomials
in the three middle families, including the empty monomial. Then
\begin{align*}
M &: \quad \{F^qXu:q\ge0,\ X\in\X\},\\
D_FM &: \quad \{F^qXu:q\in\Z,\ X\in\X\},\\
Q_c &: \quad \{F^{-m}Xu+M:m\ge1,\ X\in\X\}
\end{align*}
are bases. In particular, \(Q_c\) is nonzero and countably infinite-dimensional.
The operator \(F\) on \(Q_c\) is surjective and locally nilpotent, and
\[
\ker F^r=\Span\{F^{-m}Xu+M:1\le m\le r,\ X\in\X\}\qquad(r\ge1).
\]
Every nonzero \(\g\)-submodule of \(Q_c\) meets \(\ker F\) nontrivially.
\end{proposition}
\begin{proof}
The first basis is the induced-module PBW theorem. It makes \(M\) free
over \(\C[F]\), hence \(F\) acts injectively. Localization gives the second
basis: spanning follows by moving denominators to the left, and linear
independence follows by multiplying a putative relation by a common power
of \(F\) and using the first basis. Taking the quotient gives the third.
The formulas for \(F\) and its kernels are now immediate.
For a nonzero vector in a submodule, apply the highest power of \(F\)
whose value is still nonzero.
\end{proof}

The kernels above are vector-space filtration terms, not generally
\(\g\)-submodules. As a module over the commuting algebra
\(\C[f_j:j\le b]\), the quotient is a direct sum of copies of
\[
\bigl(\C[F,F^{-1}]/\C[F]\bigr)
\otimes\C[f_j:j\le b,\ j\ne c],
\]
indexed by the \(e,h\) PBW monomials. This is the usual first local
cohomology model for a principal ideal; it does not describe the full
\(\g\)-action by polynomial multiplication.

\begin{lemma}[Commutation and generation by pure fractions]\label{lem:fractions}
For \(s\in U\) and \(m\ge1\), one has the finite identity
\begin{equation}\label{eq:commute}
sF^{-m}=\sum_{j\ge0}\binom{m+j-1}{j}F^{-m-j}(\ad F)^j(s).
\end{equation}
Moreover, \(Q_c=\sum_{m\ge1}Uw_m\).
\end{lemma}
\begin{proof}
For \(m=1\), commute \(s\) past \(F^{-1}\) repeatedly; the process terminates
because \(\ad F\) is locally nilpotent. Induction on \(m\) gives the
binomial coefficient in \eqref{eq:commute}.

For the generation assertion, give a numerator PBW word the degree
\(\rho=2\#e+\#h\); \(f\)-factors have degree zero. Reordering does not
increase this degree, and every nonzero application of \(\ad F\) lowers
it strictly. Inducing-subalgebra factors are evaluated on \(u\) after
reordering. Thus
\[
Xw_m=F^{-m}Xu+M+
  \text{a finite sum of terms of strictly smaller numerator degree}.
\]
Induct on \(\rho(X)\), simultaneously for all denominator orders \(m\).
At degree zero the numerator contains only commuting \(f\)'s, so the
claim is immediate. Every correction term at the next degree belongs to
\(\sum_{n\ge1}Uw_n\) by induction. Subtracting the corrections yields
\(F^{-m}Xu+M\), proving the assertion for every basis vector.
\end{proof}

\begin{proposition}[The negative-root profile]\label{prop:root}
On \(Q_c\), the root modes have the following properties:
\[
\begin{array}{c|l}
\text{index}&\text{action}\\ \hline
j=c&f_j\text{ is surjective and locally nilpotent},\\
j\le b,\ j\ne c&f_j\text{ is injective and not surjective},\\
j>b&f_j\text{ is locally nilpotent}.
\end{array}
\]
Consequently, the set of locally nilpotent negative-root modes is exactly
\(\{c\}\cup\{j:j>b\}\).
\end{proposition}
\begin{proof}
The first assertion is Proposition~\ref{prop:pbw}. For the second, \(f_j\)
is multiplication by a free polynomial variable in the commuting-algebra
model. This is injective and does not hit the constant term in that variable.
For \(j>b\), use \(f_ju=0\) and local nilpotence of \(\ad f_j\) to obtain
local nilpotence on \(M\). Since \(f_j\) commutes with \(F\), it remains
locally nilpotent on \(D_FM\) and on \(Q_c\).
\end{proof}

\begin{corollary}[Nonisomorphism and Hom vanishing]\label{cor:hom}
For \(c<b\), \(Q_c\) is not isomorphic to any standard module
\(M_{a',b'}(\psi)\) of the form defined in Section~1.
For a fixed \(M\) and distinct \(c,c'\le b\),
\[
\Hom_\g(T_{f_c}M,T_{f_{c'}}M)=0.
\]
\end{corollary}
\begin{proof}
The negative-root local-nilpotence set of a standard module is the half-line
\(\{j:j>b'\}\), never \(\{c\}\cup\{j:j>b\}\) for \(c<b\).
For the Hom assertion, \(f_c\) is locally nilpotent on the source and
injective on the target. The image of every vector must therefore be zero.
This argument also covers comparison with the boundary quotient \(c'=b\).
\end{proof}

\begin{proposition}[Smoothness and weight structure]\label{prop:smooth}
The module \(Q_c\) is smooth, \(K\) acts by \(\kappa\), and \(h_0\) acts
semisimply with
\[
\Supp_{h_0}Q_c=\lambda_0+2\Z.
\]
Every nonzero weight space is countably infinite-dimensional.
The module has no nonzero finite-dimensional submodules or quotients,
and is not integrable for the full affine algebra.
\end{proposition}
\begin{proof}
If \(X\) has \(r_e\) \(e\)-factors and \(r_f\) \(f\)-factors, then
\[
h_0(F^{-m}Xu+M)
=(\lambda_0+2m+2r_e-2r_f)(F^{-m}Xu+M).
\]
All weights in the stated coset occur by varying \(m\ge1\) and the power
of any \(f_j\) with \(j\le b,\ j\ne c\).
Appending arbitrary powers of \(h_{-1}\) gives infinitely many linearly
independent vectors of each weight. PBW also gives countability.

The induced module \(M\) is smooth: commute a sufficiently high current
mode through any fixed finite PBW word until it annihilates \(u\).
Formula \eqref{eq:commute} gives
\begin{align}
h_rF^{-m}
 &=F^{-m}h_r+2mF^{-m-1}f_{r+c},\label{eq:hcomm}\\
e_rF^{-m}
 &=F^{-m}e_r-mF^{-m-1}(h_{r+c}+r\delta_{r+c,0}K)\notag\\
 &\hspace{18mm}-m(m+1)F^{-m-2}f_{r+2c}.\label{eq:ecomm}
\end{align}
For a fixed numerator, choose \(r,r+c,r+2c\) sufficiently large; all
terms vanish and the central delta is zero. The \(f_r\)'s commute with
\(F\). Hence \(D_FM\), and therefore \(Q_c\), is smooth.

Surjectivity and local nilpotence of \(F\) descend to every quotient.
On a finite-dimensional vector space these properties force the space
to be zero. A finite-dimensional smooth \(\g\)-module is trivial:
there is a common high-mode annihilation bound, and the Lie ideal
generated by those high modes is all of \(\g\), including \(K\).
But a trivial submodule of \(Q_c\) is impossible because \(f_{c+1}\)
is injective. Finally, a free root mode \(f_j\), \(j\le b,\ j\ne c\),
is not locally nilpotent, excluding integrability.
\end{proof}

Thus \(Q_c\) is a weight module, but not a finite-weight-multiplicity
Harish--Chandra module. The behavior of positive Cartan modes below
does not contradict semisimplicity of \(h_0\).

\section{Positive Cartan modes and localization depth}

\begin{proposition}[Strict growth of pole order]\label{prop:pole}
For \(1\le i\le d_*\), \(\zeta\in\C\), and \(0\ne v\in Q_c\),
\[
\pole_F((h_i-\zeta)v)=\pole_F(v)+1,
\]
where pole order means the largest denominator exponent in the PBW
expansion. In particular,
\[
\ker p(h_i)=0\qquad(0\ne p\in\C[t]).
\]
Thus \(h_i\) has no nonzero locally finite vectors, eigenvectors,
or generalized eigenvectors.
\end{proposition}
\begin{proof}
Write the top pole term of \(v\) as \(F^{-m}v_m+M\), with \(v_m\ne0\)
in the span of the numerator PBW monomials containing no \(F\).
The term \(F^{-m}h_i v_m\) in \eqref{eq:hcomm} has pole order at most \(m\).
The coefficient of pole order \(m+1\) is
\[
2m f_{c+i}v_m\pmod{FM}.
\]
Since \(c+i\le b\) and \(c+i\ne c\), multiplication by \(f_{c+i}\)
is injective on the vector-space quotient \(M/FM\).
This coefficient is therefore nonzero. Lower-pole terms and \(-\zeta v\)
cannot cancel it. Iteration shows that the vectors
\(v,h_iv,\ldots,h_i^nv\) have distinct pole orders. Hence no nonzero
polynomial in \(h_i\) annihilates \(v\).
\end{proof}

\begin{corollary}\label{cor:no-whittaker}
For every standard induced module as in Section~1,
\[
\Hom_\g(M_{a',b'}(\psi),Q_c)=0.
\]
More generally, \(Q_c\) has no nonzero character vector for a Lie
subalgebra containing \(h_1\).
\end{corollary}
\begin{proof}
Such a character vector would be an \(h_1\)-eigenvector.
A homomorphism from the induced module is determined by the image
of its generating character vector, which must therefore be zero.
\end{proof}

This is not a claim that \(Q_c\) has no Whittaker vectors for
\emph{every} possible subalgebra; the condition that the subalgebra
contain \(h_1\) is essential to this argument.

\begin{proposition}[High Cartan modes]\label{prop:high}
If \(i>d_*\), then \(h_iw_m=\lambda_iw_m\) for every \(m\ge1\),
and \(h_i-\lambda_i\) is locally nilpotent on all of \(Q_c\).
\end{proposition}
\begin{proof}
Since \(c+i>b\), the correction term in \eqref{eq:hcomm} kills \(u\).
For a general smooth module, if \(i>0\) and \(h_iw=\lambda w\), then
\[
(h_i-\lambda)^N Xw=(\ad h_i)^N(X)w.
\]
For a fixed finite word \(X\), each commutator with \(h_i\) either raises
the index of a root factor by \(i\), or replaces a Cartan factor by a
central term, which admits no further nonzero commutator.
For sufficiently large \(N\), every nonzero term contains a root factor
with arbitrarily high index. Moving it to the right can decrease its
index only by sums of the fixed negative indices in the original word.
It therefore still has sufficiently high index to annihilate \(w\).
This proves local nilpotence on \(Uw\). Apply it to all \(w_m\) and use
Lemma~\ref{lem:fractions}.
\end{proof}

Consequently, the positive Cartan modes with no nonzero locally finite
vectors are precisely \(h_1,\ldots,h_{b-c}\). This gives a second
intrinsic way to read off the localization depth.

\section{The basic module \texorpdfstring{\(Q_0=T_{f_0}M_{-1,1}\)}{Q0 = T(f0) M(-1,1)}}

Here \(a=-1\), \(b=1\), \(c=0\), and \(L=d_*=1\).
All statements in this section allow arbitrary
\(\lambda_0,\lambda_1,\kappa\). The pure fractions satisfy
\begin{align}
f_0w_1&=0,\qquad f_0w_m=w_{m-1}\quad(m>1),\notag\\
h_0w_m&=(\lambda_0+2m)w_m,\label{eq:basic}\\
e_0w_m&=-m(\lambda_0+m+1)w_{m+1},\label{eq:basic-e}\\
h_1w_m&=\lambda_1w_m+2m f_1w_{m+1},\notag\\
e_1w_m&=-m\lambda_1w_{m+1}-m(m+1)f_1w_{m+2}.\notag
\end{align}
Unlike boundary localization at \(f_1\), the pure-fraction recursion
is controlled by \(\lambda_0\), not by \(\lambda_1\).

\subsection{A canonical submodule at integral parameters}

\begin{proposition}[The \(e_0\)-locally nilpotent part]\label{prop:integer}
Let
\[
G=\{v\in Q_0:e_0^Nv=0\text{ for some }N\ge1\}.
\]
This is a \(\g\)-submodule, and
\[
\begin{cases}
0\ne G\ne Q_0,&\lambda_0\in\Z,\\
G=0,&\lambda_0\notin\Z.
\end{cases}
\]
Thus \(Q_0\) is reducible for every integral \(\lambda_0\).
For nonintegral \(\lambda_0\), \(e_0\) is injective; simplicity is not
asserted.
\end{proposition}
\begin{proof}
Set \(E=e_0\). If \(E^Nv=0\) and \((\ad E)^r(s)=0\) for all sufficiently
large \(r\), the expansion
\[
E^n(sv)=\sum_{r=0}^n\binom nr(\ad E)^r(s)E^{n-r}v
\]
vanishes for sufficiently large \(n\). Since \(\ad E\) is locally
nilpotent on \(U\), this proves \(UG\subseteq G\).

Suppose first that \(p=\lambda_0\in\Z\). Choose an integer
\(k\ge\max(1,p+2)\), and put \(x=f_0\), \(y=f_1\), \(\eta=\lambda_1\).
The subspace \(\C[x,y]u\) is stable under \(e_0,h_0,f_0\), with
\begin{equation}\label{eq:polynomial-e}
e_0(x^jy^\ell u)
=j(p-j+1-2\ell)x^{j-1}y^\ell u
 \ell\eta x^jy^{\ell-1}u.
\end{equation}
Define
\begin{equation}\label{eq:highest-vector}
v_k=\sum_{j=0}^k a_jx^jy^{k-j}u,\qquad
a_0=1,\qquad
a_j=-\frac{(k-j+1)\eta}{j(p-2k+j+1)}a_{j-1}.
\end{equation}
All denominators are nonzero, since
\(p-2k+j+1\le p-k+1<0\). In \eqref{eq:polynomial-e}, the coefficient of
\(x^{j-1}y^{k-j}u\) in \(e_0v_k\) is
\[
j(p-2k+j+1)a_j+(k-j+1)\eta a_{j-1}=0.
\]
Hence \(e_0v_k=0\), and \(h_0v_k=\nu v_k\), where \(\nu=p-2k\le-2\).
The rank-one commutation formula gives
\[
e_0(f_0^{-m}v_k)=-m(\nu+m+1)f_0^{-m-1}v_k.
\]
The coefficient vanishes at \(m=-\nu-1\). In \(Q_0\),
\[
f_0^{-1}v_k+M=f_1^kw_1\ne0,\qquad
e_0^{-\nu-1}(f_1^kw_1)=0.
\]
The first equality holds because every term of \eqref{eq:highest-vector}
with \(j\ge1\) lies in \(M\) after multiplication by \(f_0^{-1}\);
nonvanishing follows from PBW. Thus \(G\ne0\).
On the other hand, take \(m\ge\max(1,-p)\). Every coefficient in the
successive applications of \eqref{eq:basic-e} to \(w_m\) is nonzero.
Thus \(w_m\notin G\), proving that \(G\) is proper.

Now suppose \(\lambda_0\notin\Z\) and \(G\ne0\).
Because \(G\) is \(h_0\)-stable and \(h_0\) is semisimple, it contains
a nonzero weight vector. Applying a suitable power of \(e_0\) produces
a nonzero weight vector \(v'\in\ker e_0\). Write \(h_0v'=\nu v'\),
and choose the least \(N\ge1\) with \(f_0^Nv'=0\).
The \(\slc\) identity implies
\[
0=e_0f_0^Nv'=N(\nu-N+1)f_0^{N-1}v'.
\]
By minimality, \(\nu=N-1\in\Z\). This contradicts
\(\nu\in\lambda_0+2\Z\). Therefore \(G=0\).
\end{proof}

\begin{remark}[What this submodule does and does not prove]\label{rem:G}
The submodule \(G\) is the largest submodule on which \(e_0\) is locally
nilpotent. Its restriction to
\(\langle e_0,h_0,f_0\rangle\cong\slc\) is locally finite:
both root operators are locally nilpotent and \(h_0\) acts semisimply,
so the rank-one PBW theorem gives finite-dimensional cyclic submodules.
The quotient \(Q_0/G\) is nonzero and \(e_0\) acts injectively on it,
since \(e_0v\in G\) implies \(v\in G\).
Neither simplicity of \(G\) or \(Q_0/G\), nor length two or
indecomposability of \(Q_0\), has been proved here.
\end{remark}

\subsection{Cyclicity, including the exceptional parameters}

\begin{proposition}[Exact cyclicity statement]\label{prop:cyclic}
The module \(Q_0\) is cyclic for all parameters, and
\[
Q_0=Uw_1\quad\Longleftrightarrow\quad
\lambda_0\notin\{-2,-3,-4,\ldots\}.
\]
If \(\lambda_0=-s\) with \(s\ge2\), then \(Q_0=Uw_s\), while \(Uw_1\)
is a proper submodule.
\end{proposition}
\begin{proof}
If \(\lambda_0\) is not a negative integer at most \(-2\), every
coefficient in \eqref{eq:basic-e}, starting at \(m=1\), is nonzero.
Induction on \(m\) gives \(w_m\in Uw_1\) for all \(m\ge1\).
Lemma~\ref{lem:fractions}, whose induction covers every numerator and
every denominator order, then gives \(Q_0=Uw_1\).

If \(\lambda_0=-s\), the coefficient vanishes precisely at \(m=s-1\).
Starting instead at \(w_s\), repeated applications of \(e_0\) produce
all \(w_m\) with \(m\ge s\), and applications of \(f_0\) produce
\(w_{s-1},\ldots,w_1\). The same numerator-degree induction yields
\(Q_0=Uw_s\). However, \(e_0^{s-1}w_1=0\), so
\(Uw_1\subseteq G\), whereas \(w_s\notin G\).
Proposition~\ref{prop:integer} proves that \(Uw_1\ne Q_0\).
\end{proof}

Two examples distinguish cyclicity from simplicity.
If \(\lambda_0=-2\), then \(e_0w_1=f_0w_1=h_0w_1=0\), but the affine
submodule \(Uw_1\) is infinite-dimensional: the vectors \(f_1^kw_1\)
are independent. If \(\lambda_0=0\), then \(Q_0=Uw_1\), yet
\[
0\ne f_1^2w_1\in G,\qquad e_0^3f_1^2w_1=0.
\]
Thus a cyclic generator need not generate a simple module.

\subsection{An exact Takiff model}

The subspace
\[
V=\Span\{f_1^kw_m:m\ge1,\ k\ge0\}
\cong \C[x,x^{-1},y]/\C[x,y]
\]
is stable under the nonnegative current algebra \(\slc[t]\), with all
modes of degree at least two acting by zero. It is therefore a module
for the Takiff algebra \(\slc[t]/t^2\slc[t]\).
Under \(f_1^kw_m\leftrightarrow x^{-m}y^k\), the action is
\begin{align*}
f_0&=x,\qquad f_1=y,\\
h_0&=\lambda_0-2x\partial_x-2y\partial_y,\\
h_1&=\lambda_1-2y\partial_x,\\
e_0&=\lambda_0\partial_x+\lambda_1\partial_y
      -x\partial_x^2-2y\partial_x\partial_y,\\
e_1&=\lambda_1\partial_x-y\partial_x^2.
\end{align*}
These formulas follow by commuting the six modes through \(x^jy^\ell u\);
they preserve \(\C[x,y]\), hence descend to the quotient.
They are exact formulas, not a degree truncation.
The subspace \(V\) is not a submodule for the full affine algebra
(for example, negative Cartan modes leave it).
One cannot promote an arbitrary Takiff submodule to an affine submodule;
the definition of \(G\) above is what guarantees affine stability.

\section{General nonboundary indices: deep and shallow regions}

\subsection{Spectral-flow normalization}

For \(k\in\Z\), define the automorphism
\[
\sigma_k(e_i)=e_{i+k},\quad
\sigma_k(f_i)=f_{i-k},\quad
\sigma_k(h_i)=h_i+k\delta_{i,0}K,\quad
\sigma_k(K)=K.
\]
Our twist convention is \(x\cdot_{\rm new}v=\sigma_k(x)v\).
Tracking the inducing conditions and the localization element gives
\begin{equation}\label{eq:spectral}
(Q_c)^{\sigma_{-c}}\cong
T_{f_0}M_{a+c,b-c}(\chi'),
\qquad
\chi'(h_0)=\mu:=\lambda_0-c\kappa,
\end{equation}
with \(\chi'(K)=\kappa\) and \(\chi'(h_i)=\lambda_i\) for \(i>0\).
This is an isomorphism \emph{after twisting}, not an untwisted
isomorphism of the original modules.
Such automorphisms also appear in \cite[Section 3.1]{FGXZ}.
The useful normalized coordinates are
\[
A=a+c,\qquad B=b-c=d_*,\qquad A+B=L-1,\qquad \mu=\lambda_0-c\kappa.
\]

\begin{proposition}[The deep region]\label{prop:deep}
Suppose \(c\le-a-1\), equivalently \(d_*\ge L\). Put
\[
E=e_{-c},\qquad H=h_0-cK,\qquad F=f_c.
\]
These elements form an \(\slc\)-triple, \(Eu=0\), \(Hu=\mu u\), and
\[
Ew_m=-m(\mu+m+1)w_{m+1}.
\]
The following statements hold without assuming \(\lambda_L\ne0\):
\begin{enumerate}[label=(\roman*)]
\item \(Q_c=Uw_1\) if and only if \(\mu\notin\{-2,-3,\ldots\}\).
If \(\mu=-s\), \(s\ge2\), then \(Q_c=Uw_s\ne Uw_1\).
\item The subspace
\[
G_c=\{v\in Q_c:E^Nv=0\text{ for some }N\ge1\}
\]
is a \(\g\)-submodule. It is nonzero and proper when \(\mu\in\Z\),
and is zero when \(\mu\notin\Z\).
\item The operator \(E\) acts injectively on the nonzero quotient
\(Q_c/G_c\).
\end{enumerate}
\end{proposition}
\begin{proof}
The commutator \([e_{-c},f_c]=h_0-cK\) gives the rank-one triple.
The deep-region assumption is exactly \(-c>a\), so \(Eu=0\).
The formula for \(Ew_m\) follows from the rank-one commutation identity.

Here is the explicit integral-parameter construction needed to generalize
Proposition~\ref{prop:integer}; the pure-fraction recursion alone is not
sufficient for all integral parameters.
Set \(Y_b=f_b\). On \(\C[F,Y_b]u\), direct commutation gives
\begin{align}
E(F^jY_b^\ell u)
&=j(\mu-j+1-2\ell)F^{j-1}Y_b^\ell u\notag\\
&\hspace{12mm}+\ell\lambda_{d_*}F^jY_b^{\ell-1}u.\label{eq:deep-poly}
\end{align}
Indeed, \([E,Y_b]=h_{d_*}\), the \(H\)-weight of \(Y_b^\ell u\)
is \(\mu-2\ell\), and the further commutators involve \(f_{b+d_*}\),
which annihilates \(u\). Thus \eqref{eq:deep-poly} is exactly
\eqref{eq:polynomial-e} with
\((p,\eta,x,y)\) replaced by \((\mu,\lambda_{d_*},F,Y_b)\).

When \(\mu\in\Z\), choose an integer \(k\ge\max(1,\mu+2)\) and form
the highest vector \eqref{eq:highest-vector} with these replacements.
It gives
\[
0\ne f_b^kw_1\in G_c,\qquad
E^{\,2k-\mu-1}(f_b^kw_1)=0.
\]
If \(d_*>L\), then \(\lambda_{d_*}=0\), and that highest vector
simply equals \(f_b^ku\). Pure fractions with
\(m\ge\max(1,-\mu)\) are not \(E\)-locally nilpotent, so \(G_c\) is proper.
Stability follows from local nilpotence of \(\ad E\).

For \(\mu\notin\Z\), the weight support for \(H\) is \(\mu+2\Z\).
The highest-vector argument with the least \(N\) such that \(F^Nv'=0\)
in Proposition~\ref{prop:integer} again forces an integral highest weight,
a contradiction. Thus \(G_c=0\). Injectivity on \(Q_c/G_c\) follows from
\(Ev\in G_c\Rightarrow v\in G_c\).
Finally, the proof of Proposition~\ref{prop:cyclic} applies with \(\mu\)
in place of \(\lambda_0\), using Lemma~\ref{lem:fractions} to pass
from all pure fractions to the whole module.
\end{proof}

\begin{proposition}[The shallow region]\label{prop:shallow}
Suppose \(-a\le c<b\), equivalently \(1\le d_*<L\).
Then \(E=e_{-c}\) acts injectively on \(Q_c\), for every value of \(\mu\).
If \(\lambda_L\ne0\), the vector \(w_1\) generates \(Q_c\).
\end{proposition}
\begin{proof}
In this region \(e_{-c}\) belongs to the free PBW complement, so
\(Eu=0\) cannot be used. Instead set \(Y=e_{L-c}\).
Since \(L-c>a\) and \(L+c>b\), one has
\[
Yu=0,\qquad [Y,F]=h_L,\qquad f_{L+c}u=0.
\]
There is no central term in the middle identity because \(L\ge1\).
Formula \eqref{eq:ecomm} therefore yields
\[
Yw_m=-m\lambda_Lw_{m+1}.
\]
For \(\lambda_L\ne0\), induction produces every \(w_m\) from \(w_1\);
Lemma~\ref{lem:fractions} then proves cyclicity.

For injectivity of \(E\), use numerator degree \(\rho=2\#e+\#h\).
On the highest-degree part of a nonzero PBW expansion in \(Q_c\),
left multiplication by \(E\) adds a free \(e_{-c}\)-factor, increasing
degree by two. Terms from commuting \(E\) past a numerator or a
denominator have degree at most one above the original degree.
In the associated graded PBW space the highest term is multiplication
by a free polynomial variable, hence is injective. It cannot cancel.
Therefore \(E\) is injective.
\end{proof}

The deep-region integral obstruction therefore does not extend to all
indices \(c<b\). In the shallow region, whether
\(\lambda_L\ne0\) and \(\kappa\ne-2\) imply simplicity remains open here.

\section{Extensions, annihilators, and derived functors}

\begin{proposition}[The natural extension is nonsplit]\label{prop:extension}
The short exact sequence
\[
0\longrightarrow M\longrightarrow D_FM\longrightarrow Q_c\longrightarrow0
\]
does not split. In fact,
\[
\Hom_\g(Q_c,D_FM)=\Hom_\g(Q_c,M)=\Hom_\g(M,Q_c)=0.
\]
If \(M\) is simple, it is the unique simple submodule of \(D_FM\)
and an essential submodule; consequently \(D_FM\) is indecomposable.
\end{proposition}
\begin{proof}
The source \(Q_c\) is \(F\)-torsion, whereas \(F\) is injective on
\(D_FM\) and on \(M\). This proves the first two Hom vanishings and
excludes a section of the quotient map. The third is
Corollary~\ref{cor:no-whittaker}.
For any nonzero \(v\in D_FM\), a sufficiently large power of \(F\)
lies in \(M\) and is nonzero. Hence every nonzero submodule of \(D_FM\)
meets \(M\). If \(M\) is simple, this intersection equals \(M\),
so every nonzero submodule contains \(M\).
The remaining assertions follow.
\end{proof}

This proves indecomposability of \(D_FM\) when \(M\) is simple,
not indecomposability of its quotient \(Q_c\).

\begin{proposition}[Annihilator preservation]\label{prop:ann}
In the ordinary universal enveloping algebra,
\[
\Ann_U M=\Ann_U D_FM=\Ann_U Q_c.
\]
\end{proposition}
\begin{proof}
Let \(\Delta=\ad F\). The two-sided ideal \(\Ann_U M\) is
\(\Delta\)-stable. Thus \eqref{eq:commute} shows that every element of
\(\Ann_U M\) annihilates \(D_FM\). The inclusion \(M\subseteq D_FM\)
gives equality of the first two annihilators, and this common ideal
is contained in \(\Ann_U Q_c\).

For the reverse inclusion, let \(s\in\Ann_U Q_c\), and choose \(t\ge0\)
with \(\Delta^{t+1}(s)=0\). For every \(v\in M\) and \(m\ge1\),
we have \(sF^{-m}v\in M\). Multiplication by \(F^{m+t}\), using
\eqref{eq:commute}, gives
\begin{equation}\label{eq:ann-poly}
P_v(m):=\sum_{j=0}^t\binom{m+j-1}{j}
F^{t-j}\Delta^j(s)v\ \in F^{m+t}M.
\end{equation}
The left side is a vector-valued polynomial in \(m\).
Its finitely many coefficients have finite PBW support, so there is
a common upper bound on their nonnegative \(F\)-exponents.
For all sufficiently large \(m\), their span has zero intersection
with \(F^{m+t}M\). Therefore \(P_v(m)=0\) for all sufficiently large
integers \(m\), and \(P_v\) is identically zero.
Evaluating at \(m=0\), with binomial coefficients interpreted as
polynomials, yields
\[
0=P_v(0)=F^tsv,
\]
because \(\binom{j-1}{j}=0\) for \(j\ge1\).
Injectivity of \(F\) on \(M\) gives \(sv=0\). Since \(v\) was arbitrary,
\(s\in\Ann_U M\).
\end{proof}

The proposition concerns ordinary \(U\)-elements. It does not establish
an analogous statement for the completed center at the critical level.

\begin{proposition}[Functorial and derived properties]\label{prop:derived}
For an arbitrary \(U\)-module \(V\), define \(T_FV\) as the cokernel of
the localization map, without assuming that this map is injective.
Then
\begin{align*}
T_FV&=(U_F/U)\otimes_U V,\\
L_1T_F(V)&\cong\{v\in V:F^nv=0\text{ for some }n\ge0\},\\
L_jT_F(V)&=0\qquad(j\ge2).
\end{align*}
The functor is right exact and is exact on short exact sequences whose
three terms are \(F\)-torsion-free. In particular,
\[
D_FQ_c=0,\qquad T_FQ_c=0,\qquad L_1T_F(Q_c)\cong Q_c.
\]
\end{proposition}
\begin{proof}
Ore localization is flat as a right \(U\)-module. The sequence
\[
0\longrightarrow U\longrightarrow U_F\longrightarrow U_F/U
\longrightarrow0
\]
is therefore a flat resolution of the right module \(U_F/U\).
Tensoring with \(V\) identifies its degree-one homology with
\(\ker(V\to D_FV)\), which is precisely the \(F\)-torsion submodule.
There is no higher homology. The long exact sequence gives the
stated exactness on torsion-free sequences.
Finally, every vector of \(Q_c\) is \(F\)-torsion.
\end{proof}

\section{Further localization and the boundary isomorphism}

\begin{proposition}[Commuting quotient localizations]\label{prop:commuting}
Let \(F=f_c\) and \(G=f_j\) with distinct \(c,j\le b\).
Write \(D_{F,G}M\) for simultaneous localization at their powers.
There are natural isomorphisms
\[
T_G(T_FM)\cong
\frac{D_{F,G}M}{D_FM+D_GM}
\cong T_F(T_GM).
\]
The analogous statement holds for any finite family of distinct
free negative-root modes, independently of their order.
\end{proposition}
\begin{proof}
The modes commute, their localizations commute, and PBW makes \(M\)
free in the two independent polynomial variables \(F,G\).
Inside \(D_{F,G}M\), it follows that
\[
D_FM\cap D_GM=M.
\]
Moreover, \(G\) is injective on \(T_FM\), since it still acts as a
free polynomial variable. Localizing the exact sequence defining
\(T_FM\) at \(G\) gives
\[
D_G(T_FM)\cong D_{F,G}M/D_GM.
\]
Taking the quotient by the image of \(T_FM\) gives the middle
expression in the proposition. Its PBW basis consists exactly of
terms with both exponents strictly negative. Interchanging \(F,G\)
gives the other isomorphism. The same argument works for finitely
many variables.
\end{proof}

Recall the boundary result established in the earlier notes \cite{Boundary}:
if \(\lambda_L\ne0\), then
\begin{equation}\label{eq:boundary}
T_{f_b}M_{a,b}(\chi)\cong M_{a+1,b-1}(\chi^{(1)}),
\qquad
\chi^{(1)}(h_0)=\lambda_0+2,
\end{equation}
with the other character values unchanged. The present note uses this
previously proved result; it does not replace its detailed isomorphism proof.

\begin{corollary}[Filling the interval of missing modes]\label{cor:fill}
Assume \(\lambda_L\ne0\). Applying one quotient localization at each of
\[
f_c,f_{c+1},\ldots,f_b
\]
in any order yields
\[
M_{a+b-c+1,c-1}\bigl(\chi^{(b-c+1)}\bigr),
\qquad
\chi^{(b-c+1)}(h_0)=\lambda_0+2(b-c+1),
\]
and all other character values are unchanged.
\end{corollary}
\begin{proof}
By Proposition~\ref{prop:commuting}, arrange the operations in the
order \(f_b,f_{b-1},\ldots,f_c\). At each step the next mode is a
boundary mode for the resulting standard module.
Apply \eqref{eq:boundary} repeatedly. The value of \(a+b+1=L\)
and the nonzero parameter \(\lambda_L\) are unchanged throughout.
\end{proof}

For the basic example this reads
\begin{align*}
T_{f_1}\bigl(T_{f_0}M_{-1,1}(\chi)\bigr)
&\cong T_{f_0}\bigl(T_{f_1}M_{-1,1}(\chi)\bigr)\\
&\cong M_{1,-1}(\chi^{(2)}),\qquad \lambda_1\ne0.
\end{align*}
These are iterated \emph{quotient} localizations. They are not the
quotient of simultaneous localization by \(M\) alone: one must divide
by the sum of the relevant partial localizations.
Even if the final module is simple, simplicity of the first quotient
does not follow.

\section{Adjoining the derivation and the vertex-algebra viewpoint}

Let
\[
\widetilde{\g}=\g\rtimes\C d,\qquad
[d,x_r]=r x_r,\qquad [d,K]=0,
\]
where \(x=e,f,h\). Put \(\widetilde U=U(\widetilde{\g})\) and
\(\I(V)=\widetilde U\otimes_U V\). PBW gives
\[
\I(V)\cong\C[z]\otimes V,
\]
with actions
\begin{align*}
d(P(z)\otimes v)&=zP(z)\otimes v,\\
x_r(P(z)\otimes v)&=P(z-r)\otimes x_rv,\\
K(P(z)\otimes v)&=P(z)\otimes Kv.
\end{align*}
In particular, adjoining \(d\) is not the naive tensor-product action
with unchanged current modes.

\begin{proposition}[Compatibility with derivation induction]\label{prop:d}
There is a natural \(\widetilde{\g}\)-module isomorphism
\[
\I(Q_c)\cong T_{f_c}\I(M).
\]
If \(0\ne N\ne Q_c\) is a \(\g\)-submodule, then
\[
0\ne\I(N)\ne\I(Q_c).
\]
\end{proposition}
\begin{proof}
The canonical map \(\I(M)\to\I(D_FM)\) has an \(F\)-invertible target:
\[
F(P(z)\otimes v)=P(z-c)\otimes Fv,\qquad
F^{-1}(P(z)\otimes v)=P(z+c)\otimes F^{-1}v.
\]
The universal property of localization extends it uniquely to a
\(\widetilde{\g}\)-module map
\[
\alpha:D_F\I(M)\longrightarrow\I(D_FM).
\]
Its explicit formula is
\begin{equation}\label{eq:alpha}
\alpha\bigl(F^{-m}(P(z)\otimes v)\bigr)
=P(z+mc)\otimes F^{-m}v.
\end{equation}
The inverse sends
\(P(z)\otimes F^{-m}v\) to \(F^{-m}(P(z-mc)\otimes v)\).
These formulas respect equivalent fractions: replacing
\(F^{-m}v\) by \(F^{-m-1}Fv\) is precisely compensated by the
shift in the action of \(F\). Thus they define inverse maps.
The map \(\alpha\) identifies the embedded copy of \(\I(M)\)
on each side.

The algebra \(\widetilde U\) is free as a right \(U\)-module, so
\(\I\) is exact. Taking quotients therefore identifies the target
quotient with \(\I(D_FM/M)=\I(Q_c)\), as required.
Finally, \(\I\) is faithful because its underlying vector space is
\(\C[z]\otimes V\). Applying exactness and faithfulness to
\(0\to N\to Q_c\to Q_c/N\to0\) proves the last assertion.
\end{proof}

Thus the proper submodules \(G_c\) in the deep integral region remain
proper and nonzero after adjoining \(d\). No general preservation
of simplicity in the opposite situation is claimed.

By smoothness and the fixed action of \(K\), each \(Q_c\) is a weak
module for the universal affine vertex algebra \(V^\kappa(\slc)\);
see \cite[Section 2]{FGXZ}. Descending to its simple quotient
requires a separate check that the defining vertex ideal annihilates
the module. Smoothness alone does not imply this.

\section{Open problems and limits of the present results}

The following are questions, not theorems of this note.

\begin{enumerate}[leftmargin=*,label=\arabic*.]
\item \textbf{Nonintegral simplicity in the basic case.}
Under \(\lambda_1\ne0\) and \(\kappa\ne-2\), the natural candidate is
\[
T_{f_0}M_{-1,1}(\chi)\text{ simple}
\quad\Longleftrightarrow\quad \lambda_0\notin\Z.
\]
Only the necessary direction is proved here. The missing step is
to start with an arbitrary nonzero submodule and reduce one of its
vectors to a pure fraction that generates the whole quotient.
Cyclicity of \(w_1\) and injectivity of \(e_0\) do not establish this.

\item \textbf{The corresponding deep-region criterion.}
Under \(\lambda_L\ne0\), \(\kappa\ne-2\), and \(d_*\ge L\),
does \(\mu=\lambda_0-c\kappa\notin\Z\) imply simplicity?
A rank-one twisting or completion argument, or a direct PBW reduction,
may be useful. Its applicability to these smooth modules must be proved;
one cannot simply invoke a statement confined to BGG category
\(\mathcal O\).

\item \textbf{The integral-parameter submodule structure.}
Determine whether \(G_c\) and \(Q_c/G_c\) are simple, whether \(G_c\)
is the socle, whether their extension splits, and whether \(Q_c\)
has finite Jordan--Hölder length. At present, neither length two
nor indecomposability is asserted.

\item \textbf{Simplicity in the shallow region.}
Here \(e_{-c}\) is injective for all \(\mu\), so the deep-region
integral obstruction disappears. Does simplicity of the original
module suffice, or are there other resonance conditions?

\item \textbf{Parameter isomorphisms and degenerations.}
The level \(\kappa\), support \(\lambda_0+2\Z\), negative-root
local-nilpotence set, and positive-Cartan local-finiteness profile
give necessary invariants. For \(i>d_*\), the generalized eigenvalue
\(\lambda_i\) is an additional invariant.
A complete parameter classification, especially for \(\lambda_L=0\)
or \(\kappa=-2\), remains unresolved.

\item \textbf{Critical-level central reduction.}
At \(\kappa=-2\), study the completed center, its central reductions,
their compatibility with localization, and the resulting simple
quotients. Proposition~\ref{prop:ann} does not settle completed-center
questions.

\item \textbf{Multiple missing modes and twisted localization.}
Proposition~\ref{prop:commuting} suggests a family indexed by finite
sets of free modes. One may also study generalized conjugation
\[
\Theta_\xi(s)=\sum_{j\ge0}\binom{\xi}{j}(\ad F)^j(s)F^{-j},
\qquad \xi\in\C,
\]
on \(D_FM\), where the sum is finite for each \(s\in U\).
This construction requires \(F\) to be invertible; it is not
directly an operation on \(Q_c\).
\end{enumerate}

\clearpage
\section{Sources and verification record}

The original induced modules, their simplicity criterion, spectral-flow
automorphisms, and the smooth-module vertex-algebra correspondence are
drawn from \cite{FGXZ}. The localization results of \cite{FGM} provide
relevant comparisons: Theorem~4.3 exhibits an integral rank-one
reducibility condition \(J-iK\in\Z\) for imaginary Verma/free-field
modules, and Section~5 constructs primitive vectors and proves
the corresponding nonintegral simplicity statements.
Those modules differ from the present ones, so their sufficiency
theorems are not used as proofs of our open simplicity questions.

The boundary isomorphism and the derivation-induction framework were
established in the earlier project notes \cite{Boundary}.
This document is an English organization of the current nonboundary
results, not a revision of those earlier boundary proofs.

The accompanying computation
\texttt{work/check\_nonboundary\_takiff.py} recorded 5,866 exact
rational checks of the six Takiff generators' bracket relations,
the highest-vector recurrence at integral parameters, explicit
locally nilpotent vectors, and pure-fraction resonance.
The checks include \(\lambda_1=0\) and do not truncate polynomial
degree in the action formulas. They nevertheless test only finitely
many parameters and vectors and do not replace the general proofs.
No new Danus run was used for these nonboundary results.
No independent audit or Lean/Coq formal verification is claimed.

\begingroup
\raggedright
\begin{thebibliography}{9}
\bibitem{FGXZ}
V.~Futorny, X.~Guo, Y.~Xue, and K.~Zhao,
\emph{Smooth representations of affine Kac--Moody algebras},
arXiv:2404.03855v2 (2024; journal version 2025).
\href{https://arxiv.org/html/2404.03855v2}{arXiv full text}.

\bibitem{FGM}
V.~Futorny, D.~Grantcharov, and R.~A.~Martins,
\emph{Localization of free field realizations of affine Lie algebras},
arXiv:1404.7148v2 (2015).
\href{https://arxiv.org/html/1404.7148v2}{arXiv full text}.

\bibitem{Boundary}
Earlier project working notes, \emph{Boundary localization: advisor note}
(September 2026) and \emph{Boundary localization and derivation extension}
(September 8, 2026). English source files:
\texttt{boundary\_localization\_advisor\_note\_en.tex} and
\texttt{boundary\_derivation\_extension\_en.tex}.
\end{thebibliography}
\endgroup
\end{document}
